\section{模型的评价与推广}

\subsection{模型的优点}

\textbf{其一，四问共用一套可核验的守恒离散框架。}本文没有为每个问题单独编写求解器，
而是把控制方程骨架、边界处理与守恒格式统一起来，只在物性口径、时程与坐标口径上分支。
这样既避免了同一物理过程被重复推导，也使四问的结果具有直接可比性。

\textbf{其二，质量守恒可以在离散层面被逐时步核验。}顶点有限体积的通量形式使全域水量的变化
与表面通量之间存在精确的离散恒等式，残差达到机器精度；对移动边界问题，本文进一步给出
“关闭传质后收缩不改变干基含水率”的解析不变量检验，使表观对流项的正确性可以被\textbf{证伪}，
而不是仅凭收敛曲线判断。

\textbf{其三，主导不确定性被识别并量化。}分析表明烘干时长对题目给定的传质系数并不敏感，
而主要由模型口径决定。这一结论改变了结论的呈现方式：论文不应给出一个孤立的烘干时长数字，
而应同时给出它所依赖的口径与相应的影响幅度。

\textbf{其四，报告的数值可逐条溯源。}所有进入论文的数值都来自运行记录，并由机器生成的宏引用，
没有一个是手工抄录的；表格亦由脚本从运行结果直接生成，从根本上排除了抄写误差。

\subsection{模型的缺点与改进方向}

\textbf{其一，维度简化。}模型忽略了药材两端面的轴向输运。在 $30$~min 的预热窗口内这在数量级上成立，
但在数日的干燥时程内轴向渗透可达数厘米每端，必须作为局限声明。若要改进，可把模型扩展为二维轴对称问题，
或在端面处补充轴向通量的等效边界条件。

\textbf{其二，质量方程的读法。}本文按题目明文的 Fick 标准形式建式，密度不进入质量方程，
因此几何收缩在模型内不是质量汇。替代读法会使烘干时长变化近 $15\%$，该差异只能通过口径声明而非数值手段消除。

\textbf{其三，经验式的适用性。}物性完全采用题目给定的经验式，未作独立验证；
扩散系数的强非线性使结果对经验式的形式较为敏感，这是模型本身无法回避的不确定性。

\textbf{其四，判据尾部的平坦性。}干燥判据位于含水率曲线的平坦段，使烘干时长对数值噪声敏感。
改进方向是改用\textbf{时间区间}而非单个时刻来报告烘干完成，或同时给出判据附近的曲线形态。

\subsection{模型的推广}

本文的建模框架不限于中药材。凡是满足“圆柱（或球形）几何、内部扩散控制、表面为对流边界”的干燥、
吸附或浸提过程，都可以直接套用：把物性经验式替换为目标材料的数据，把环境序列替换为实际的工艺曲线，
并保留本文的守恒检验与不变量检验即可。问题四的归一化坐标方法尤其具有推广价值——
凡计算域随时间单调变化（收缩、膨胀、溶解、烧蚀）且变化规律已知的扩散问题，都可以用同一手法
把移动边界转化为已知的表观对流项，从而在固定网格上求解并保留精确的守恒性。
